\begin{prob}
    \begin{subprobset}
        \begin{subprob}
            Recursive code is not necessarily more efficient than iterative code.
            In fact, it is easy to write bad recursive code that is very slow
            if we aren't careful. Here's an example:

            \inputminted{python}{\thisdir/include/part_a.py}

            This code might look fine, but it's actually needlessly slow. The reason
            is that \python{numbers[1:]} creates a \emph{copy} of the list,
            and a copy takes time linear in the size of the new list (since that much
            memory has to be allocated).

            Write down a recurrence relation for \python{summation} and solve it
            to determine the function's asymptotic time complexity.

            \begin{soln}
                The copy takes time $\Theta(n)$, and the recursive call is on an array
                of size $n - 1$. Therefore our recurrence is:
                \[
                    T(n)
                    =
                    \begin{cases}
                        T(n-1) + \Theta(n),& n = \geq 1\\
                        \Theta(1), n = 0 
                    \end{cases}
                \]
                We'll replace the $\Theta$ notation to solve the recurrence: $T(n) = T(n-1) + n$.
                Unrolling a few times, we find:
                \begin{align*}
                    T(n) 
        &=
        T(n-1) + n
        \\
        &=
        T(n-2) + (n - 1) + n
        \\
        &=
        T(n-3) + (n-2) + (n - 1) + n
        \\
                \end{align*}
                In general, it looks like $T(n) = T(n - k) + n + (n-1) + \ldots + (n - k + 1)$.
                The base case is reached when $k = n$, and on that iteration we have:
                \[
                    T(n) = T(n - n) + n + (n-1) + \ldots + (n - n + 1) = \Theta(n^2)
                \]

                So this algorithm takes quadratic time complexity!
            \end{soln}

        \end{subprob}

        \begin{subprob}
            Taking a long time to run is not the only way that an algorithm can be inefficient;
            it can also take a lot of memory. Just as we used time complexity to measure the
            speed of an algorithm, we can use \emph{space complexity} to measure how much
            memory a piece of code requires.
            The space complexity of an algorithm is the amount of
            \emph{auxiliary} (or \emph{extra}) memory it requires to execute.

            It turns out that \python{summation} is not only slow; it is wasteful
            of memory, too. When \python{numbers[1:]} is executed in the root call,
            a copy is made, resulting in the creation of a new list containing $n - 1$
            elements. This list is kept in memory while all subsequent recursive calls are made.

            At the moment the base case is reached, what is the total size of all lists
            created by copying in all recursive calls made so far?
            State your answer exactly (and not in asymptotic notation), simplifying
            it as much as possible.

            \begin{soln}
                The root-level call on an array of size $n$ creates a list with $n - 1$ elements. The first recursive call,
                which is on an array of size $n-1$, creates a list with $n - 2$ elements. So
                the $k$th call (where $k = 1$ is the root call) is on an array of size $n - k +1$, and creates
                a list of size $n - k$.

                The base case occurs when the input list is empty; that is, when $n - k + 1 = 0$. Solving for
                $k$, we find $k = n + 1$. But no copy is made on this call (it returns before a copy is made).
                So the total size of all created lists is:
                \[
                    \sum_{k=1}^{n} (n - k) = n^2 - \frac{n(n+1)}{2} = \frac{n(n-1)}{2}
                \]
            \end{soln}
        \end{subprob}

    \end{subprobset}
\end{prob}
